\section{Tipo de investigación}

Por el tipo de problema, el trabajo cabe en tres casilleros:

\begin{itemize}
  \item \textbf{Investigación tecnológica.} Partimos de un problema de ingeniería. Los artefactos
        son el contrato y el \ServiceDomain{}; el producto es un procedimiento que se puede evaluar.
  \item \textbf{Investigación basada en diseño (\textit{design science}).} El modelo toma contratos
        OpenAPI bancarios y \ServiceDomain{}s BIAN y devuelve S, E, C y el \AlignmentScore{}.
  \item \textbf{Investigación no experimental:} no emplea población ni muestra humana; no aplica
        encuestas ni instrumentos de recolección de percepciones.
\end{itemize}

La validez del estudio se sustentará en la reproducibilidad del procedimiento, la trazabilidad de
las métricas y el análisis de amenazas a la validez (Capítulo~\ref{cap:metodologia}), más que en
inferencia estadística sobre personas.
